\begin{textsample}{POS Dim 3 – pi – Score 29.00 – pi/pi\_inf\_6.txt}  \label{ex:f3_pos_251}
5.
ORGANIZAÇÃO DA EDUCAÇÃO \textbf{INFANTIL} NO DOCUMENTO CURRICULAR A intencionalidade da educação \textbf{infantil} consiste na organização e proposição, pelo educador, de experiências que permitam às \textbf{crianças} conhecer a si e ao outro e de conhecer e compreender as relações com a natureza, com a cultura e com a produção científica.
Desta forma, o currículo da educação \textbf{infantil} do Estado do Piauí está organizado em campos de experiência, objetivos de aprendizagem e período de transição para o ensino fundamental.
Na imagem a seguir, apresenta -se esta organização : A \textbf{criança} piauiense possui diversas características que as identificam e as diferenciam entre si.
Negros, brancos, amarelos, pardos, surdos ou ouvintes, indígenas, quilombolas, com “ lócus ” diferenciados vivem na cidade ou no campo, no litoral, nas comunidades, bairros e cada uma delas detentora de direitos, de acesso à educação.
A BNCC elege três faixas \textbf{etárias} que reordenam de forma pedagógica e \textbf{lúdica} as ações integradas de escola e família a saber : \textbf{bebês}, \textbf{crianças} bem \textbf{pequenas} e \textbf{crianças} \textbf{pequenas}.
Desta forma o documento soma -se às realidades, necessitando de atualização em projetos políticos pedagógicos, plano estadual e municipal de educação.
Este documento inicia uma nova trajetória e representa um marco na garantia de direitos das \textbf{crianças} do Piauí e do Brasil.
Na educação \textbf{infantil}, as aprendizagens e o desenvolvimento das \textbf{crianças} têm como eixos estruturantes as interações e as \textbf{brincadeiras} e, como direitos de aprendizagens : conviver, \textbf{brincar}, participar, explorar, expressar -se e conhecer -se.
As \textbf{crianças} têm \textbf{vontades} e interesses, necessitam de um ambiente de convivência onde elas possam se expressar e vivenciar suas experiências e se conhecer melhor.
No quadro a seguir, explica -se como o currículo da educação \textbf{infantil} do Piauí está estruturado e conectado com os ideais de \textbf{criança} e de educação.
Quadro I – Direitos de Aprendizagem da Educação \textbf{Infantil} | Direitos de Aprendizagem | Ação de Vivências | |---|---| | Conviver | Com outras \textbf{crianças} e \textbf{adultos}, em \textbf{pequenos} e grandes grupos, utilizando diferentes linguagens, ampliando o conhecimento de si e do outro, o respeito em relação à cultura e às diferenças entre as pessoas. | | \textbf{Brincar} | Cotidianamente de diversas formas, em diferentes espaços e tempos, com \textbf{crianças} e \textbf{adultos}, ampliando e diversificando seu acesso a produções culturais, seus conhecimentos, sua imaginação, sua criatividade, suas experiências emocionais, corporais, \textbf{sensoriais}, expressivas, cognitivas, sociais e \textbf{relacionais}. | | Participar | Ativamente, com \textbf{adultos} e outras \textbf{crianças}, tanto do planejamento da gestão da escola e das atividades propostas pelo educador quanto da realização das atividades da vida cotidiana, tais como a escolha das \textbf{brincadeiras}, dos materiais e dos ambientes, desenvolvendo diferentes linguagens e elaborando conhecimentos, decidindo e se posicionando. | | Explorar | Movimentos, \textbf{gestos}, \textbf{sons}, formas, \textbf{texturas}, cores, palavras, emoções, transformações, relacionamentos, histórias, objetos, elementos da natureza, na escola e fora dela, ampliando seus saberes sobre a cultura, em suas diversas modalidades : as artes, a escrita, a ciência e a tecnologia. | | Expressar | Como sujeito dialógico, criativo e sensível, suas necessidades, emoções, sentimentos, dúvidas, hipóteses, descobertas, opiniões, questionamentos, por meio de diferentes linguagens. | | Conhecer -se | E construir sua identidade pessoal, social e cultural, constituindo uma imagem positiva de si e de seus grupos de pertencimento, nas diversas experiências de \textbf{cuidados}, interações, \textbf{brincadeiras} e linguagens vivenciadas na \textbf{instituição} escolar e em seu contexto familiar e comunitário. | Fonte : ( BRASIL 2017, p. 36 ) Dessa forma, a educação \textbf{infantil} tem como foco um padrão de qualidade e garantia de direitos que se viabilizam pela articulação das competências gerais com os campos de experiências, direitos de aprendizagens e eixos norteadores, definidas pela Base Nacional Comum Curricular.
De acordo a Lei nº 12.796, de 2013, Art. 31, a educação \textbf{infantil} será organizada seguindo os seguintes princípios : I – avaliação mediante \textbf{acompanhamento} e \textbf{registro} do desenvolvimento das \textbf{crianças}, sem o objetivo de promoção, mesmo para o acesso ao ensino fundamental ; II – carga horária \textbf{mínima} anual de 800 ( oitocentas ) horas, distribuída por um mínimo de 200 ( duzentos ) \textbf{dias} de trabalho educacional ; III – atendimento à \textbf{criança} de, no mínimo, 4 ( quatro ) horas \textbf{diárias} para o turno parcial e de 7 ( sete ) horas para a jornada integral ; IV – controle de frequência pela \textbf{instituição} de educação pré-escolar, exigida a frequência \textbf{mínima} de 60 \% ( sessenta por cento ) do total de horas ; V – expedição de documentação que permita atestar os processos de desenvolvimento e aprendizagem da \textbf{criança}.
Nesse sentido, a BNCC define que a educação \textbf{infantil} será organizada atendendo as faixas \textbf{etárias} de 0 a 1 ano e 6 \textbf{meses} de idade ; de 1 ano e 7 \textbf{meses} a 3 anos e 11 \textbf{meses} e de 4 a 5 anos de idade.
A BNCC organiza a Educação \textbf{Infantil} por grupos \textbf{etários} e não mais por \textbf{creche} e \textbf{pré-escola}.
A \textbf{pré-escola} compreende a faixa \textbf{etária} de \textbf{crianças} \textbf{pequenas} de 4 a 5 anos e 11 \textbf{meses}.
A \textbf{creche} passa a ser identificada pelos \textbf{bebês} com idade de zero a 1 ano e 6 \textbf{meses} de idade e também pelas \textbf{crianças} bem \textbf{pequenas} com 1 ano e 7 \textbf{meses} a 3 anos e 11 \textbf{meses}.
Os municípios do Estado do Piauí não possuem estrutura suficiente para atender os três grupos de faixa \textbf{etária} organizado pela BNCC, por isso devem adequar os espaços de atendimento à educação \textbf{infantil}.
A educação \textbf{infantil} precisa de uma estrutura para atender as \textbf{crianças} de acordo com a organização dos níveis e faixa \textbf{etária}. - \textbf{Bebês} de zero a 1 ano e 6 \textbf{meses} : essa \textbf{criança} tem seus ritmos próprios, necessitam de espaços para engatinhar, rolar, ensaiar os primeiros passos, explorar materiais diversos, observar, \textbf{brincar}, tocar o outro, alimentar -se, tomar banho, repousar, dormir, \textbf{satisfazendo}, assim, suas necessidades essenciais.
Dessa forma, recomenda -se que o espaço para essa idade esteja situado em local silencioso, preservado das áreas de grande \textbf{movimentação} e proporcione \textbf{conforto} às \textbf{crianças}.
A estrutura do espaço para os \textbf{bebês} devem dispor de : salas para as atividades, salas para repouso, fraldário, lactário para a higienização e solário para o banho de sol. - \textbf{Crianças} bem \textbf{pequenas} de 1 ano e 7 \textbf{meses} a 3 anos e 11 \textbf{meses} e as \textbf{Crianças} \textbf{pequenas} de 4 a 5 anos e 11 \textbf{meses} necessitam de um espaço físico visto como um ambiente que possibilita e contribui para a vivência e a expressão das culturas \textbf{infantis} como os jogos, \textbf{brincadeiras}, músicas, histórias que expressam a especificidade da \textbf{criança}.
Por isso deve -se organizar um ambiente adequado à proposta pedagógica da \textbf{instituição}, que possibilite à \textbf{criança} a realização de explorações, interações e \textbf{brincadeiras}, garantindo -lhe identidade, segurança, confiança, interações socioeducativas e privacidade, promovendo oportunidades de aprendizagem e desenvolvimento integral.
A \textbf{rotina} das \textbf{instituições} de educação \textbf{infantil} traz atividades planejadas pelos professores que lidam com o espaço e o tempo a todo o momento.
Dessa forma, as \textbf{instituições} de ensino devem \textbf{apoiar} os professores na organização desses espaços, organização do tempo de \textbf{brincar}, de tomar banho, de se alimentar e de repousar.
Organizar o cotidiano das \textbf{crianças} da Educação \textbf{Infantil} pressupõe pensar que o estabelecimento de uma sequência básica de atividades \textbf{diárias} é, antes de mais nada, o resultado da leitura que fazemos do nosso grupo de \textbf{crianças}, a partir, principalmente, de suas necessidades.
É importante que o educador observe o que as \textbf{crianças} \textbf{brincam}, como estas \textbf{brincadeiras} se desenvolvem, o que mais gostam de fazer, em que espaços preferem ficar, o que lhes chama mais atenção, em que momentos do \textbf{dia} estão mais tranquilos ou mais agitados.
Este conhecimento é fundamental para que a estruturação espaço-temporal tenha significado.
Ao lado disto, também é importante considerar o contexto sociocultural no qual se insere e a proposta pedagógica da \textbf{instituição}, que deverão lhe dar suporte ( BARBOSA ; HORN, 2001, p. 67 ).
Dessa forma, o olhar dos professores para organizar esses espaços e o tempo das atividades da educação \textbf{infantil} precisam da atenção dos gestores porque eles são muito importantes no sentido de proporcionar as condições necessárias para o desenvolvimento da prática pedagógica na etapa da educação \textbf{infantil}.
A sistemática de avaliação na educação \textbf{infantil} é processual, por isso ela acontece no \textbf{dia} a \textbf{dia}, durante o período de aprendizado e desenvolvimento da \textbf{criança}.
Para isso, as escolas de educação \textbf{infantil} devem criar procedimentos para o \textbf{acompanhamento} do trabalho pedagógico e para a avaliação do desenvolvimento das \textbf{crianças}, sem objetivo de reprovação, seleção, promoção ou classificação.
Devem garantir a observação das atividades, das \textbf{brincadeiras} e interações das \textbf{crianças} no cotidiano, utilizando -se de \textbf{registros} realizados pelos professores como relatórios, fotografias, desenhos e álbuns.
Os professores observarão a continuidade dos processos de aprendizagens por meio da criação de estratégias adequadas aos diferentes momentos vividos pelas \textbf{crianças}.
As escolas de educação \textbf{infantil} terão documentação específica que permita às famílias conhecer o trabalho da \textbf{instituição} com as \textbf{crianças} e os processos de desenvolvimento e aprendizagem da \textbf{criança} na educação \textbf{infantil}.
Nessa etapa de educação as \textbf{crianças} não ficam retidas ( Res CNE/CEB nº 5/2009, art. 10 ).
Dessa forma, a avaliação na educação \textbf{infantil} se dá principalmente pela observação sistemática, \textbf{registro} em portfólio ou caderno de anotações e relatório.

% matched lemmas: acompanhamento, adulto, apoiar, bebê, brincadeira, brincar, conforto, creche, criança, cuidado, dia, diário, etário, gesto, infantil, instituição, lúdico, movimentação, mês, pequeno, pré-escola, registro, relacional, rotina, satisfazer, sensorial, som, textura, vontade
\end{textsample}
